% GENERATED FILE. Edit canonical lessons, then run npm run generate:books.
% canonical-source-hash: fnv1a64:a4eadb4c3f330997
% canonical-lessons: HI-C527-kala, HI-C527-citra, HI-C527-kuci, HI-C527-rekha, HI-C527-rangna

\chapter{Colours on the Wall}
\label{ch:hi-colours-on-the-wall}
\hlchaptermodality{527}

\begin{chapteropening}
\textbf{By the end of this chapter:} \emph{I can talk about pictures, colours and lines.}

Complete the last lesson of chapter 527: I can talk about pictures, colours and lines.
\end{chapteropening}

\section[\texorpdfstring{\hi{कला} (kalā)}{कला (kalā)}]{\hi{कला} (kalā) — art}

\label{lesson:HI-C527-kala}



\begin{quote}
Before the new one: say the Hindi for an ace (in cards), then the Hindi for chess.
\end{quote}

\subsection*{You'll want to know: \hi{कला}}
\textbf{\hi{कला}} — \emph{kalā} — \textquotedblleft{}art\textquotedblright{}.

\begin{quote}
\hi{हमारे} \hi{स्कूल} \hi{में} \hi{हर} \hi{बुधवार} \hi{कला} \hi{की} \hi{कक्षा} \hi{होती} \hi{है।}
\end{quote}

\begin{quote}
\emph{Hamāre skūl meṁ har budhvār kalā kī kakṣā hotī hai.}
\end{quote}

\begin{quote}
Our school has an art class every Wednesday.
\end{quote}

Feminine.

\begin{cousinweb}
\textbf{\hi{कला}} (\emph{kalā}) is a tatsama, taken whole from Sanskrit \emph{kalā}, \textquotedblleft{}a small part; any practical art\textquotedblright{} (Monier-Williams).
\end{cousinweb}

\begin{grammarlens}[title={What you've built}]
One more word for describing what you lived and felt.
\end{grammarlens}

\subsection*{Your turn}
\begin{itemize}
\raggedright
  \item \emph{Say it:} \emph{kalā}
  \item \emph{Say it:} \emph{kalā}, once more
  \item \emph{Say it:} \emph{Hamāre skūl meṁ har budhvār kalā kī kakṣā hotī hai.} — Our school has an art class every Wednesday.
  \item \emph{Recall:} say the Hindi for an ace (in cards), then the Hindi for chess, then say \emph{kalā} again
\end{itemize}

\begin{tcolorbox}[breakable,colback=teal!4,colframe=teal!35!black,title={Before you move on}]
What does \hi{कला} mean? (\textquotedblleft{}Art\textquotedblright{}.) Say it once more.
\end{tcolorbox}
\section[\texorpdfstring{\hi{चित्र} (citra)}{चित्र (citra)}]{\hi{चित्र} (citra) — a picture, a painting}

\label{lesson:HI-C527-citra}



\begin{quote}
Before the new one: say the Hindi for chess, then the Hindi for art.
\end{quote}

\subsection*{You'll want to know: \hi{चित्र}}
\textbf{\hi{चित्र}} — \emph{citra} — \textquotedblleft{}a picture, a painting\textquotedblright{}.

\begin{quote}
\hi{दीवार} \hi{पर} \hi{एक} \hi{सुंदर} \hi{चित्र} \hi{है।}
\end{quote}

\begin{quote}
\emph{Dīvār par ek sundar citra hai.}
\end{quote}

\begin{quote}
There is a beautiful painting on the wall.
\end{quote}

Masculine.

\begin{cousinweb}
\textbf{\hi{चित्र}} (\emph{citra}) is a tatsama, taken whole from Sanskrit \emph{citra}, \textquotedblleft{}bright, bright-coloured; a picture\textquotedblright{} (Monier-Williams).
\end{cousinweb}

\begin{grammarlens}[title={What you've built}]
One more word for describing what you lived and felt.
\end{grammarlens}

\subsection*{Your turn}
\begin{itemize}
\raggedright
  \item \emph{Say it:} \emph{citra}
  \item \emph{Say it:} \emph{citra}, once more
  \item \emph{Say it:} \emph{Dīvār par ek sundar citra hai.} — There is a beautiful painting on the wall.
  \item \emph{Recall:} say the Hindi for chess, then the Hindi for art, then say \emph{citra} again
\end{itemize}

\begin{tcolorbox}[breakable,colback=teal!4,colframe=teal!35!black,title={Before you move on}]
What does \hi{चित्र} mean? (\textquotedblleft{}A picture, a painting\textquotedblright{}.) Say it once more.
\end{tcolorbox}
\section[\texorpdfstring{\hi{कूची} (kūcī)}{कूची (kūcī)}]{\hi{कूची} (kūcī) — a paintbrush}

\label{lesson:HI-C527-kuci}



\begin{quote}
Before the new one: say the Hindi for art, then the Hindi for a picture, a painting.
\end{quote}

\subsection*{You'll want to know: \hi{कूची}}
\textbf{\hi{कूची}} — \emph{kūcī} — \textquotedblleft{}a paintbrush\textquotedblright{}.

\begin{quote}
\hi{चित्रकार} \hi{ने} \hi{कूची} \hi{से} \hi{आसमान} \hi{नीला} \hi{किया।}
\end{quote}

\begin{quote}
\emph{Citrakār ne kūcī se āsmān nīlā kiyā.}
\end{quote}

\begin{quote}
The painter made the sky blue with a brush.
\end{quote}

Feminine.

\begin{cousinweb}
\textbf{\hi{कूची}} (\emph{kūcī}) is a tadbhava, inherited from Sanskrit \emph{kūrca}, \textquotedblleft{}a bunch, a brush\textquotedblright{} (CDIAL 3408).
\end{cousinweb}

\begin{grammarlens}[title={What you've built}]
One more word for describing what you lived and felt.
\end{grammarlens}

\subsection*{Your turn}
\begin{itemize}
\raggedright
  \item \emph{Say it:} \emph{kūcī}
  \item \emph{Say it:} \emph{kūcī}, once more
  \item \emph{Say it:} \emph{Citrakār ne kūcī se āsmān nīlā kiyā.} — The painter made the sky blue with a brush.
  \item \emph{Recall:} say the Hindi for art, then the Hindi for a picture, a painting, then say \emph{kūcī} again
\end{itemize}

\begin{tcolorbox}[breakable,colback=teal!4,colframe=teal!35!black,title={Before you move on}]
What does \hi{कूची} mean? (\textquotedblleft{}A paintbrush\textquotedblright{}.) Say it once more.
\end{tcolorbox}
\section[\texorpdfstring{\hi{रेखा} (rekhā)}{रेखा (rekhā)}]{\hi{रेखा} (rekhā) — a line}

\label{lesson:HI-C527-rekha}



\begin{quote}
Before the new one: say the Hindi for a picture, a painting, then the Hindi for a paintbrush.
\end{quote}

\subsection*{You'll want to know: \hi{रेखा}}
\textbf{\hi{रेखा}} — \emph{rekhā} — \textquotedblleft{}a line\textquotedblright{}.

\begin{quote}
\hi{उसने} \hi{काग़ज़} \hi{पर} \hi{एक} \hi{सीधी} \hi{रेखा} \hi{खींची।}
\end{quote}

\begin{quote}
\emph{Usne kāġaz par ek sīdhī rekhā khīñcī.}
\end{quote}

\begin{quote}
He drew a straight line on the paper.
\end{quote}

Feminine.

\begin{cousinweb}
\textbf{\hi{रेखा}} (\emph{rekhā}) is a tatsama, taken whole from Sanskrit \emph{rekhā}, \textquotedblleft{}a line, a stroke\textquotedblright{} (Monier-Williams).
\end{cousinweb}

\begin{grammarlens}[title={What you've built}]
One more word for describing what you lived and felt.
\end{grammarlens}

\subsection*{Your turn}
\begin{itemize}
\raggedright
  \item \emph{Say it:} \emph{rekhā}
  \item \emph{Say it:} \emph{rekhā}, once more
  \item \emph{Say it:} \emph{Usne kāġaz par ek sīdhī rekhā khīñcī.} — He drew a straight line on the paper.
  \item \emph{Recall:} say the Hindi for a picture, a painting, then the Hindi for a paintbrush, then say \emph{rekhā} again
\end{itemize}

\begin{tcolorbox}[breakable,colback=teal!4,colframe=teal!35!black,title={Before you move on}]
What does \hi{रेखा} mean? (\textquotedblleft{}A line\textquotedblright{}.) Say it once more.
\end{tcolorbox}
\section[\texorpdfstring{\hi{रंगना} (raṅgnā)}{रंगना (raṅgnā)}]{\hi{रंगना} (raṅgnā) — to colour; to paint; to dye}

\label{lesson:HI-C527-rangna}



\begin{quote}
Before the new one: say the Hindi for a paintbrush, then the Hindi for a line (drawn).
\end{quote}

\subsection*{You'll want to know: \hi{रंगना}}
\textbf{\hi{रंगना}} — \emph{raṅgnā} — \textquotedblleft{}to colour; to paint; to dye\textquotedblright{}.

\begin{quote}
\hi{बच्चों} \hi{ने} \hi{दीवार} \hi{हरे} \hi{रंग} \hi{से} \hi{रंगी।}
\end{quote}

\begin{quote}
\emph{Baccoṁ ne dīvār hare raṅg se raṅgī.}
\end{quote}

\begin{quote}
The children painted the wall green.
\end{quote}

Past, with \hi{ने} (he or she …): \hi{उसने} \hi{कुर्सी} \hi{रंगी}, \emph{usne kursī raṅgī}.

\begin{cousinweb}
\textbf{\hi{रंगना}} (\emph{raṅgnā}) is the verb of \emph{raṅg}, \textquotedblleft{}colour\textquotedblright{}, probably from the same source as Hindi \emph{rãgānā}, \textquotedblleft{}to be coloured\textquotedblright{}, which CDIAL files under \emph{*raṅgayati}, \textquotedblleft{}dyes\textquotedblright{}, a reconstructed form (CDIAL 10570).
\end{cousinweb}

\begin{grammarlens}[title={What you've built}]
One more word for describing what you lived and felt.
\end{grammarlens}

\subsection*{Your turn}
\begin{itemize}
\raggedright
  \item \emph{Say it:} \emph{raṅgnā}
  \item \emph{Say it:} \emph{raṅgnā}, once more
  \item \emph{Say it:} \emph{Baccoṁ ne dīvār hare raṅg se raṅgī.} — The children painted the wall green.
  \item \emph{Recall:} say the Hindi for a paintbrush, then the Hindi for a line (drawn), then say \emph{raṅgnā} again
\end{itemize}

\begin{tcolorbox}[breakable,colback=teal!4,colframe=teal!35!black,title={Before you move on}]
What does \hi{रंगना} mean? (\textquotedblleft{}To colour; to paint; to dye\textquotedblright{}.) Say it once more.
\end{tcolorbox}
